% ****** Start of file apssamp.tex ******
%
%   This file is part of the APS files in the REVTeX 4.2 distribution.
%   Version 4.2a of REVTeX, December 2014
%
%   Copyright (c) 2014 The American Physical Society.
%
%   See the REVTeX 4 README file for restrictions and more information.
%
% TeX'ing this file requires that you have AMS-LaTeX 2.0 installed
% as well as the rest of the prerequisites for REVTeX 4.2
%
% See the REVTeX 4 README file
% It also requires running BibTeX. The commands are as follows:
%
%  1)  latex apssamp.tex
%  2)  bibtex apssamp
%  3)  latex apssamp.tex
%  4)  latex apssamp.tex

%This sample document demonstrates proper use of REV\TeX~4.2 (and \LaTeXe) in mansucripts prepared for submission to APS journals. Further information can be found in the REV\TeX~4.2 documentation included in the distribution or available at \url{http://journals.aps.org/revtex/}.

%When commands are referred to in this example file, they are always shown with their required arguments, using normal \TeX{} format. In this format, \verb+#1+, \verb+#2+, etc. stand for required author-supplied arguments to commands. For example, in \verb+\section{#1}+ the \verb+#1+ stands for the title text of the author's section heading, and in \verb+\title{#1}+ the \verb+#1+ stands for the title text of the paper.

%Line breaks in section headings at all levels can be introduced using \textbackslash\textbackslash. A blank input line tells \TeX\ that the paragraph has ended. Note that top-level section headings are automatically uppercased. If a specific letter or word should appear in lowercase instead, you must escape it using \verb+\lowercase{#1}+ as in the word ``via'' above.
%
\documentclass[%
 reprint, %This file may be formatted in either the \texttt{preprint} or \texttt{reprint} style. \texttt{reprint} format mimics final journal output.  Either format may be used for submission purposes. \texttt{letter} sized paper should be used when submitting to APS journals.
%superscriptaddress,
%groupedaddress,
%unsortedaddress,
%runinaddress,
%frontmatterverbose, 
%preprint,
%preprintnumbers,
%nofootinbib,
%nobibnotes,
%bibnotes,
amsmath,amssymb,
aps,
%pra,
%prb,
%rmp,
%prstab,
%prstper,
%floatfix,
]{revtex4-2}
\usepackage{float}
\usepackage{graphicx}% Include figure files
\usepackage{dcolumn}% Align table columns on decimal point
\usepackage{bm}% bold math
\usepackage{hyperref}% add hypertext capabilities
%\usepackage[mathlines]{lineno}% Enable numbering of text and display math
%\linenumbers\relax % Commence numbering lines

%\usepackage[showframe,%Uncomment any one of the following lines to test 
%%scale=0.7, marginratio={1:1, 2:3}, ignoreall,% default settings
%%text={7in,10in},centering,
%%margin=1.5in,
%%total={6.5in,8.75in}, top=1.2in, left=0.9in, includefoot,
%%height=10in,a5paper,hmargin={3cm,0.8in},
%]{geometry}
\begin{document}
\preprint{APS/123-QED}


\title{Emergent Tetragonality in a Fundamentally Orthorhombic Material}% Force line breaks with \\


\author{Anisha G. Singh$^{1,2,3}$}
\author{Maja D. Bachmann$^{1,2,3}$}
\author{Joshua J. Sanchez$^4$}
\author{Akshat Pandey$^5$}
\author{Aharon Kapitulnik$^{1,2,3,5}$}
\author{Jong Woo Kim$^6$}
\author{Philip J. Ryan$^6$}
\author{Steven A. Kivelson$^{1,2,5,7}$}
\author{Ian R. Fisher$^{1,2,3}$}
\begin{abstract}
[1] Geballe Laboratory for Advanced Materials, Stanford University, Stanford, CA 94305

[2] Stanford Institute for Materials and Energy Sciences, SLAC, 2575 Sand Hill Road, Menlo 

Park, CA 94025

[3] Department of Applied Physics, Stanford University,  Stanford, CA 94305

[4] Massachusetts Institute of Technology, Cambridge, MA 02139

[5] Department of Physics, Stanford University,  Stanford, CA 94305

[6] Advanced Photon Source, Argonne National Lab, Lemont, IL 60439

[7] Rudolf Peierls Centre for Theoretical Physics, University of Oxford, Oxford, UK

\end{abstract}


\maketitle
\date{\today}% It is always \today, today,
             %  but any date may be explicitly specified

%\keywords{Suggested keywords}%Use showkeys class option if keyword
                              %display desired

%\tableofcontents
\textbf{Symmetry plays a key role in determining the physical properties of materials. By Neumann's principle, the properties of a material are invariant under the symmetry operations of the space group to which the material belongs. Continuous phase transitions are associated with a spontaneous reduction in symmetry. (For example, the onset of ferromagnetism spontaneously breaks time reversal symmetry.) Much less common are examples where proximity to a continuous phase transition leads to an increase in symmetry. Here, we find an emergent tetragonal symmetry close to an apparent charge density wave (CDW) bicritical point in a fundamentally orthorhombic material, ErTe$_3$, for which the CDW phase transitions are tuned via anisotropic strain. The underlying structure of the material remains orthorhombic for all applied strains, including at the  bicritical point, due to a glide plane symmetry in the crystal structure. Nevertheless, the observation of a divergence in the anisotropy of the in-plane elastoresistivity reveals an emergent electronic tetragonality near the  bicritical point.}

\section{\label{sec:level1}Introduction\protect\\} %The line break was forced 

ErTe$_3$, like other rare earth tritellurides (RTe$_3$), belongs to the orthorhombic space group Cmcm ($D^{17}_{2h}$, no 63). The structure comprises bilayers of almost square-planar Te nets, separated along the $b$-axis by RTe block layers. At room temperature, the $a$ and $c$ lattice parameters are almost identical, $a=0.999c$, and it is possible to tune these values using externally applied anisotropic strain, such that the material is tuned from $a>c$ to $a<c$. Despite this near equivalence of the in-plane lattice parameters, the presence of a glide plane between the Te bilayers (illustrated in Figure \ref{fig1}(a)) makes the material fundamentally orthorhombic. The glide plane operation is a non-symmorphic symmetry element and cannot be removed by external strains. Thus, even when the system is strain-tuned to a point where the in-plane lattice parameters are equivalent ($a=c$), the system does not possess a 4-fold rotational symmetry, and hence can never be truly tetragonal. 

%The ability to restore 4-fold rotational symmetry with uniaxial strain in the rare earth tritellurides, is highly nontrivial because of the symmetry elements of the Cmcm (D172h, no 63) space group to which they belong. This space group contains a glide plane element which is the origin of the system's orthorhombicity.\ref{fig1}a To understand this structure, it is instructive to first consider the simpler rare earth ditelluride, RTe$_2$. RTe$_2$, a tetragonal system, is composed of square tellurium nets which are separated by corrugated RTe block layers. In contrast, RTe$_3$ contains the same structural motifs but a different stacking pattern, built by layering two RTe$_2$ units with a reflection and a horizontal translation at 45 degrees with respect to the Te-Te bond in the Te layers. This operation is the glide plane operation which is a non-symmorphic symmetry element and cannot be removed by external strains. Since strain cannot remove the glide plane, even when the system is strain-tuned to a point where the in-plane lattice parameters are equivalent, the system is still not tetragonal. 



%\section{Strain Tuning Charge Density Wave Transitions to a Bicritical Point}


The consequences of strain-tuning become more interesting when considering its effect on ErTe$_3$'s charge density wave (CDW) state. RTe$_3$ crystals host a unidirectional, incommensurate CDW. For un-strained ErTe$_3$, this onsets via a continuous phase transition at a critical temperature $T_{CDW}$ = 268K, with the wavevector oriented along the $c$-axis.\cite{Ru2006}\cite{Ru_th} The instability is apparently driven, at least in large part, by a strongly q-dependent electron-phonon coupling, and softening of the associated phonon mode which has been observed via inelastic x-ray scattering \cite{Maschek2015} \cite{Maschek2018}. The structural motifs that make up the building blocks of RTe$_3$ result in a near-equivalence of the electronic structure and phonon spectrum in the two inequivalent in-plane directions. This near-equivalence is reflected in a substantial softening of the phonon mode at the CDW wavevector in the transverse ($a$-axis) direction, even while the phonon in the $c$-axis direction approaches zero energy as $T \to T_{CDW}$.  Earlier calculations for a simplified model indicate that unidirectional order is favored for sufficiently strong electron-phonon coupling even on a truly square lattice \cite{Yao2006}. Consequently, the inequivalence of the $a$ and $c$ axes in ErTe$_3$ has generally been thought to act as a weak symmetry breaking field, favoring the $c$-axis CDW state over the competing $a$-axis state. %The observation of a subsequent phase transition at an even lower temperature to a 2q state comprising simultaneous $c$-axis and $a$-axis CDWs confirms the ability of the material to sustain a CDW in the a-axis direction (ref Rob Moore), while the variation of both as a function of choice of rare earth clearly reveals these to be competing phases \cite{Ru2008}.

The presence of CDW fluctuations along both the $a$ and $c$-axis directions \cite{Maschek2015} as well as the subsequent $a$-axis transition in ErTe$_3$ \cite{Ru2008} suggest the possibility to tune between these two CDW states, with lattice strain being an ideal tuning parameter. The CDW order parameters $\phi_a$ and $\phi_c$ couple to normal strains ($\epsilon_{xx},\epsilon_{yy}$ and $\epsilon_{zz}$) as $\lambda^{jj}_i \epsilon_{jj}|\phi_i|^2$ where $\lambda_i^{jj}$ are coupling constants. There are two immediate consequences of this coupling. First, a spontaneous strain (i.e. distortion) occurs as the material is cooled through the CDW transition (i.e. a nonzero value of $\phi_c$ or $\phi_a$ will result in a non-zero value of $\epsilon_{xx},\epsilon_{yy}$ and $\epsilon_{zz}$) - this is indeed observed, with the lattice growing longer in the direction of the CDW wavevector, and shorter in the transverse direction \cite{Ru2008}. Second, T$_{CDW}$ can be tuned by strain \cite{Straquadine2020}.

No symmetry operations of the point group relate the coefficients $\lambda_c^{xx}$ and $\lambda_a^{xx}$. However, the near equivalence of the electronic structure and phonon spectrum in these two directions implies that the coefficients might be closely related (i.e. it is not unreasonable to anticipate that $\lambda_c^{xx} \sim \lambda_a^{yy} \sim -\lambda_a^{xx}$ etc). If this is indeed the case, then relatively small strains would be able to not only tune $T_{CDW}$, but reverse the direction of the primary CDW wave-vector (\ref{fig1}(c)). Previous elastoresistance and elastocaloric effect measurements point towards this possibility, though direct evidence has thus far been lacking \cite{Straquadine2020}. 

\begin{figure}
\includegraphics[scale=0.5]{fig1.jpg}
\caption{\label{fig1} (a) Crystal structure of ErTe$_3$, which comprises Te bilayers separated by ErTe blocks. Note that in the Cmcm space group setting the in-plane lattice parameters are $a$ and $c$, while the long $b$ axis is perpendicular to the Te planes. A glide plane between the Te bilayers (illustrated) makes the material fundamentally orthorhombic. (b) Schematic diagram illustrating the neck of the titanium 'bowtie' platform used for applying strain to a thin single crystal of ErTe$_3$ which is bonded to its surface. Insets indicate coordinate axes for the experiment (xyz) and the crystal axes of the sample (abc). Labels indicate the effective length $L_x$ and width $L_z$ of the platform. (c) Schematic diagram illustrating the proposed phase diagram of ErTe$_3$ as a function of $\Delta L_x/L_x$. Double solid lines represent continuous phase transitions while dashed lines indicate first order phase transitions. Insets indicate deformation of the Te lattice as a result of the strain. The putative bicritical point is labeled.}
\end{figure}

It is important to note that while both $\phi_a$ and $\phi_c$ are superficially similar unidirectional CDW states, from a symmetry perspective they are in fact different ordered states, precisely because the crystal structure is different in these two directions. Consequently, properties of the two CDW states are also necessarily different. At a minimum, this means it is not given that strain should necessarily be able to rotate the CDW direction, and hence this must be established experimentally. Here, this is determined by performing high resolution x-ray diffraction experiments for samples held under variable anisotropic strain conditions. As shown in Fig. 1c, the strain-tuned phase boundaries of the two distinct CDW states meet. within experimental uncertainty, at a  bicritical  point below which there is a first order transition between the two ordered states.  Because the experiment is carried out at fixed strain, the first order transition results in a region of two-phase coexistence - i.e. domains of the two distinct CDW states - for intermediate values of the macroscopic strain $\Delta L_x/L_x$ and temperatures below the bicritical temperature. Noting that there are no symmetry-imposed reasons dictating the existence of such a multicritical point, it is intriguing to find signatures in elastoresistivity measurements of an emergent tetragonal symmetry near the strain-tuned bicritical point. 



%Significantly, tuning the two CDW phases such that their phase boundaries cross would imply a strain-tuned bicritical point, as shown in figure 1(c). 

%Though there are no symmetry-imposed reasons dictating the existence of such a multicritical point, nevertheless an emergent tetragonal symmetry arises near this strain-tuned bicritical point in this fundamentally orthorhombic material.  

%Here we report a series of both high-resolution x-ray diffraction and transport measurements under uniaxial strain to demonstrate the ability of the CDW wavevector to indeed be reoriented with strain and the appearance of a diverging elastoresistivity response at a critical strain state consistent with the emergence of a higher symmetry state proximate to a strain tuned bicritical state.



\section{High Resolution X-ray Diffraction}

Utilizing high resolution x-ray diffraction (XRD) at a synchrotron light source, the evolution of CDW superlattice peaks with externally induced strain can be observed, providing a direct probe of the behavior of the CDW state. XRD measurements were completed at the Advanced Photon Source at Argonne National Lab at Sector 6-ID-B. Here, a unique sample environment has been developed to study samples with XRD with in-situ strain tuning at cryogenic temperatures.\cite{Sanchez2019} To apply variable strain to the samples a CS-130 Razorbill strain cell was utilized.\cite{Hicks2014}  

RTe$_3$ are very micaceous crystals. To maintain the crystallinity of the sample necessary for high resolution XRD while strain-tuning, rather than mounting the sample directly into the Razorbill cell, the sample was affixed to a 'bowtie' platform designed for high strain transmission (Figure \ref{fig1}(b) \cite{Park2020}). Various sample configurations were tested; here we focus on a geometry for which tensile strain is applied along the crystallographic $a$-axis. Displacement is measured in the $x$ direction by a capacitive sensor, yielding an estimate for $\Delta L_x/L_x$. For thin samples, strain in the $z$ direction is determined by the Poisson ratio of the Ti platform. The sample itself must be sufficiently thin to mitigate strain relaxation over its  height. Finite element simulations indicated that the strain over the height of the sample is approximately homogeneous for crystals less than 50um thick. RTe$_3$ crystals exfoliate easily due to weak van der Waals bonding between the tellurium sheets, making it straightforward to cleave samples to thin dimensions. Similarly, a thin and stiff glue layer is necessary to mitigate strain relaxation. For sample mounting, Angstrom Bond epoxy was used which yielded glue layers between 5-10um thick. Samples used for XRD measurement where between 10-20um thick, while samples used for transport measurements were between 1-10um thick. In this range of thicknesses, all samples behaved similarly with strain, however thinner samples were found to have worse crystalinity and hence were avoided for XRD measurement. 

The strain-dependence of two representative CDW superlattice peaks are shown in Figure 2(a), confirming that anisotropic strain indeed reorients the primary CDW wavevector. In the absence of strain, the superlattice peak is observed along the L direction consistent with $c$-axis CDW order. As tension is applied along the $a$-axis, effectively reversing the anisotropy of the system, the L axis peak is suppressed while the H axis peak emerges, consistent with reorientation of the CDW wavevector by 90 degrees. This measurement also reveals a range of intermediate strain states where both the $c$ and $a$-axis CDW are simultaneously observed. 

In figure 2(b), the integrated intensity of both superlattice peaks is plotted as a function of $\Delta L_x/L_x$. The total CDW intensity is constant with strain, but the relative weight shifts between the $c$ and $a$ axis CDW states. 

In Figure 2(c), the behavior of a representative structural Bragg peak with an in-plane component, (2,20,0), is plotted for the same strain range, reported here as the $a$-axis lattice parameter. For the strains at which the mixed CDW phase is observed (using approximately 10 and 90 percent criteria), the Bragg peak splits, indicative of the emergence of domains. In this strain range, by rotating the CDW wavevector, the system adopts structural domains through which it can internally relax the applied stress. If, however, the applied strain exceeds the internal strain created by CDW rotation and consequent domain formation, a mono-domain state is achieved, in which the lattice parameter varies monotonically with strain. Significantly, the observation of such structural domains provides clear evidence that the simultaneous observation of $q_a$ and $q_c$ order is not the consequence of the emergence of a uniform homogeneous state comprising both wavevectors. Moreover, no peak was observed at the wavevector $q_a+q_c$.

\begin{figure} 
\includegraphics[scale=0.44]{xray_figure.jpg}
\caption{\label{fig2}(a) Observed superlattice peak intensities at (1, 27, q$_{CDW}$) and (q$_{CDW}$, 27, 1) at 200K for different values of $\Delta L_x/L_x$. Under tension, the $c$-axis CDW peak decreases in intensity while the $a$-axis CDW peak increases in intensity. Significant broadening is observed in these structural peaks due to sample quality. (b) Integrated intensity of CDW superlattice peaks shown in (a) after background subtraction and normalization to adjacent Bragg peak intensities. Shaded region indicates mixed $q_c$, $q_a$ phase (c) 2D plot of the intensity of the (2,20,0) Bragg peak as a function of strain at 200K. For intermediate strains, two peaks appear in the data due to the presence of domains, corresponding to regions of the material for which the CDW wavevector is oriented along the $a$ and $c$ axes respectively. Here, the lattice parameters of the system are unchanged as a function of strain, since the system accommodates the applied strain by changing the relative population of the two types of domain. } 
\end{figure}

\section{Elastoresistivity}

%Measurements of elastoresistivity, or the changes of the samples resistivity with strain, are highly sensitive to changes in symmetry. Generally we can define a nematic order parameter, $\Phi$, which is proportional to $\frac{\rho_{aa} - \rho_{cc}}{\rho_{aa} + \rho_{cc}}$, the in-plane resistivity anisotropy. The associated nematic susceptibility, $\chi$, would then be defined as $\frac{\delta\Phi}{\delta(\epsilon_{aa}-\epsilon_{cc})}|_{\epsilon\rightarrow 0}$, or $\frac{\delta\frac{\rho_{aa} - \rho_{cc}}{\rho_{aa} + \rho_{cc}}}{\delta(\epsilon_{aa}-\epsilon_{cc})}|_{\epsilon\rightarrow 0}$, the anisotropy of the in-plane elastoresistivity. In the case of RTe$_3$, since the system always has distinct a and c axes, we do not expect a divergence in the nematic susceptibility. If a strain state is found however where the elastoresistivity response of the system is maximized this would indicate the system has been tuned to a point where the a and c directions are electronically degenerate.

Transport measurements can provide a useful window on electronic behavior close to phase transitions, with important information encoded in its dependence on both temperature and strain respectively. Here, we focus on the in-plane resistivity anisotropy $(\rho_a-\rho_c)$, which is sensitive to the orientation of the CDW wavevector.

To perform accurate measurements of the in-plane resistive and elastoresistive anisotropies, a novel transport device was designed. Finite element simulations were performed to establish optimal sample geometry and dimensions to achieve large and homogeneous strain transmission, as well as equipotentials that are closely parallel, necessary for resistivity measurement (see Figure \ref{fig3}(a)). The device design incorporates a transverse measurement geometry which allows for direct study of the resistive anisotropy in a single measurement \cite{Walmsley2017}\cite{Shapiro2016}. The sample geometry was defined using a focused ion beam (FIB; Figure \ref{fig3}(b)) and utilizing the same design of titanium platform that was used for the x-ray diffraction measurements. In order to perform transport measurements, the platform was initially heated to high temperatures for several hours to create an insulating oxide layer prior to sample bonding. Additional features of the device design are illustrated in Figure \ref{fig3}. 

Measured values of the in-plane resistive anisotropy are plotted in Figure \ref{fig3}(c). For temperatures above the onset of CDW order, $T> T_{CDW}$, $\rho_a-\rho_c$ is close to, but not exactly zero, as expected since the two axes, though nearly equivalent, are fundamentally distinct. Below the CDW transition, opening of the CDW gap results in an increase in the resistivity, though how this affects $\rho_a$ and $\rho_c$ depends on the orientation of the CDW wavevector. When, the wave vector orients along the c-axis, there is a larger change in $\rho_a$ compared to $\rho_c$, resulting in a positive value of $\rho_a-\rho_c$ \cite{Walmsley2017}. This behavior has been previously understood in terms of a simple Boltzmann transport approach including anisotropy in the Fermi velocity.\cite{Monceau2014} Conversely, when the CDW orients along the a-axis, the change in $\rho_c$ is now larger than that of $\rho_a$ resulting in a negative value of $\rho_a-\rho_c$ \cite{Straquadine2020}.

\begin{figure}
\includegraphics[scale=0.79]{fig3.jpg}
\caption{\label{fig3} (a) Finite element simulation, utilizing COMSOL Multiphysics, of the strain transmission and current flow for a typical transport device. Gray lines show cuts in the ErTe$_3$ crystal, which in practice are made using a FIB. To achieve high and homogeneous strain transmission, the sample length along the strain direction should be maximized, resulting in the choice of a diamond shape for the active area, rather than more traditional bar shaped sample geometry. For this wider sample, multiple current injection points must then be defined in order to facilitate homogeneous current flow in the sample. (b) SEM photo of a typical device measured. Signal measured between the V$_T + / -$ contacts, the transverse voltage, correspond to the in-plane resistive anisotropy $(\rho_{a} – \rho_{c})/2$ (c) In-plane resistive anisotropy as a function of temperature for different (representative) values of compressive and tensile strain applied parallel to the $a$-axis of the sample.} 
\end{figure}

The resistivity data shown in Figure 3(c) also reveal a change in $T_{CDW}$ with strain, with higher values of $T_{CDW}$ under tension (blue curves). Taking the peak in the temperature derivative of the longitudinal resistivity $(\rho_a+\rho_c)/2$ as a measure of the critical temperature (see Appendix B), $T_{CDW}$ can be determined as a function of strain, and is shown in Figure 4(a). This data reveals a minimum value of $T_{CDW}$ close to a critical value of $\Delta L_{x}/L_x\sim 0.13\%$. Notably, this is the same strain range for which $\rho_a-\rho_c$ does not increase or decrease below $T_{CDW}$ (Figure 3(c)), implying that $T_{CDW}$ at this strain value marks the putative bicritical point at which the onset of $\phi_a$ and $\phi_c$ CDW order coincides. As discussed earlier, there is no symmetry relation between the $a$ and $c$-axis CDW states. Consequently the strain derivative of the critical temperature of the two states is not compelled to be the same, and this is indeed found to be the case (i.e. $T_{CDW}$ is not an even function about the critical strain in Figure 4(a)).

Inspection of the data shown in Figure 3(c) also reveals that the largest changes in resistivity with respect to strain occur when the resistive anisotropy is smallest (i.e. for strains closest to the critical value). To further investigate this, elastoresistance measurements were performed as a function of strain and temperature close to $T_{CDW}$. Utilizing an AC strain technique, in which a small AC strain is superimposed on an offset DC strain and the resulting change in resistivity is measured using a demodulation technique \cite{Hristov2018}, we are able to directly and sensitively measure changes in the strain derivative of the resistive anisotropy. This value, normalized by the longitudinal resistivity, is plotted as a function of the relative temperature $(T-T_{CDW})$ in Figure 4(b), for strains below, at and above the critical strain. In the temperature range above $T_{CDW}$, the material is uniform (no domains), so the microscopic strain $\epsilon_{xx}$ is equal to the measured macroscopic deformation $\Delta L_{x}/L_{x}$. Similarly, $\epsilon_{zz}$ can be estimated based on the Poisson ratio of the Ti platform. Systematic sources of uncertainty associated with strain relaxation in glue layers and strain transmission to the active area of the bow-tie device lead to systematic shifts in values, but no change in the qualitative behavior. To determine the relative temperature, $T_{CDW}$ as a function of $\Delta L_{x}/L_x$ was first determined by fitting the data shown in Figure 4(a) to a smooth polynomial. The elastoresistivity clearly peaks upon approaching $T_{CDW}$ at the critical strain value (red data), but has a smaller magnitude and rolls over for all other values of applied strain (grey data).

%As previously discussed, it is anticipated and observed that $T_{CDW}$ is highly tunable with strain as shown in Figure 4a. T$_{CDW}$ increases for both positive and negative strains at a growing rate and is minimized near $\Delta L_{x}/L_{x} = 0.13\%$. In figure 4b, the elastoresistivity is plotted at fixed strain, $\Delta L_{x}/L_{x}$ as a function of relative temperature $T-T_{CDW}(\Delta L_{xx}/L_{xx})$. Above $T_{CDW}$, a non-zero, and growing elastoresistivity is measured consistent with what is expected for growing, strain tunable CDW fluctuations in this temperature region. Moreover, this measurement confirms that the value of the elastoresistivity is highly dependent on how close to $T_{CDW}$ the measurement is taken. 

The same elastoresistivity coefficient is plotted as a function of strain in Figure 4(c) for different values of the relative temperature. This reveals a clear maximum centered around the critical strain 0.13$\%$. This is a remarkably robust observation and is not sensitive to the exact functional form used for the strain dependence of $T_{CDW}$ (see Appendix C). 

\begin{figure}
\includegraphics[scale=0.453]{fig4.jpg}
\caption{\label{fig4} (a) The strain dependence of T$_{CDW}$, as determined from the peak in the temperature derivative of the longitudinal resistivity $(\rho_{a}+\rho_{c})/2$ (see Appendix B), and expressed here as a function of $\Delta L_x/L_x$. Vertical band marks the critical strain at which the strain-derivative of $T_{CDW}$ changes sign and its uncertainty. (b) In-plane elastoresistivity anisotropy versus temperature for fixed offset strains (labelled in legend). For strain values close to the critical strain (red curve) this quantity diverges as T approaches $T_{CDW}$. While the strains in this plot are reported as a constant, this data was taken for a fixed voltage applied to the piezoelectric stack. The strain label is therefore an average over the temperature sweep with an uncertainty of +/-0.02$\%$. For this reason, the strain at which a maximum response is observed in these temperature sweeps differs slightly from the maximum strain response reported in panel (c), where instead strain sweeps at constant temperature are reported. (c) In-plane elastoresistive anisotropy, $\eta$ as defined in Section 4, versus strain for fixed relative temperature. The minimum value of $\eta$ for a given relative temperature is subtracted from all data at that temperature to facilitate comparison across temperatures on a single plot. Vertical band marks the critical strain, and its uncertainty.} 
\end{figure}

%Therefore to make a consistent comparison of the elastoresistivity measurements taken at different strains, rather than comparing data for fixed values of absolute temperature, data is compared for fixed values of relative temperature, T - $T_{CDW}(\Delta L_{x}/L_{x})$. These are the values of elastoresistivity plotted in figure 4c which can be compared to investigate if a maximum response emerges at a critical strain. This measurement clearly shows that there is indeed an intermediate strain value for which the anisotropy of the in-plane elastoresistivity is maximized consistent with the expectation for tuning to a tetragonal symmetry. This strain value is similar to the strain for which we see a minimum of $T_{CDW}$ and consistent to where we see a disappearance of the resistive anisotropy through the CDW transition. The highlighted elastoresistivity response plotted in Figure 4b is that measured at this critical strain point. Having tuned the system to a tetragonal symmetry at this strain, the measured growing elastoresistivity is indicative of nematic fluctuations.

\section{Discussion}
The differential ``nematic'' elastoresistance
\begin{align}
    \eta \equiv \frac 1 {(\rho_a+\rho_c)} \ \frac {\partial (\rho_a-\rho_c)}{\partial (\epsilon_{xx}-\epsilon_{zz})}
    \label{eta}
\end{align}
has no special significance in an orthorhombic material. The three normal strains, $\epsilon_{xx}$, $\epsilon_{yy}$ and $\epsilon_{zz}$ all belong to A$_{1g}$ representations in an orthorhombic point group, such that each strain can separately affect each of the terms $\rho_a$, $\rho_b$ and $\rho_c$ in the resistivity tensor. Consequently, while this quantity is allowed to vary with strain and temperature, there are no symmetry constraints or requirements that would dictate singular behavior as a function of temperature or strain for this particular combination of coefficients. 

In contrast, for a tetragonal material with a two-component order parameter ($\phi_x,\phi_y$) the same elastoresistance coefficient (rotating the coordinate axes such that x and y directions are related by a 4-fold symmetry) has a very specific physical significance: In this case a nematic order parameter can be defined: $N = <\phi_x^2> - <\phi_y^2>$, a quantity which, above the critical temperature, measures differences in the fluctuations of the two components of the order parameter. Since $N$ couples linearly to antisymmetric strain $\epsilon_{B1g}=(\epsilon_{xx}-\epsilon_{yy})/2$, one can define a nematic susceptibility for the same symmetry channel, $\chi_{B_{1g}} = \frac{\partial N}{\partial \epsilon_{B_{1g}}}$. In this instance, as has been extensively discussed in the context of Fe-based superconductors,\cite{Chu2012}\cite{Shapiro2015} \cite{Kuo2016} the resistive anisotropy is linearly proportional to $N$ for small values of $N$, and hence the elastoresistivity coefficient $\eta =\partial [(\rho_{xx}-\rho_{yy})/(\rho_{xx}+\rho_{yy})]/\partial \epsilon_{B_{1g}}$ is proportional to $\chi_{B_{1g}}$.\cite{Shapiro2015} Hence, a divergence of this elastoresistivity coefficient in a tetragonal material is evidence of a diverging nematic susceptibility. Such behavior has been observed for a variety of materials for which multicomponent order parameters are relevant, incluing Fe-based superconductors and BaNi$_2$As$_2$.\cite{Chu2012}\cite{Eckberg202}    

The above behavior is not, however, a consequence of symmetry for an orthorhombic material, like ErTe$_3$, for which $\phi_a$ and $\phi_c$ do not belong to a multicomponent order parameter, and for which a ‘permanent’ inequivalence of x and z directions means that $N$ has a finite value for all temperatures, and all strains. Such a system would not show any singular behavior in the elastoresistivity beyond standard Fisher-Langer-like anomalies at the strain-tuned phase transition.\cite{Hristov2019} The divergence of $\eta$ for ErTe$_3$ upon approaching the apparent bicritical point is therefore anomalous, and indeed is suggestive of an emergent tetragonal symmetry proximate to this special point in the temperature-strain plane.   

The data obtained thus far, while highly suggestive, do not permit a critical scaling analysis approaching infinitesimally close to the proposed bicritical point. Stringent requirements of, for example, strain homogeneity, compositional uniformity and temperature stability, render this a challenging task. Given these practical limitations, we initially adopt a simple mean field model to provide insight to this remarkable observation. 

To quartic order in the two order parameters, and neglecting gradient terms and the elastic energy, the contribution to the free energy density from the CDW is given by
%
\begin{equation} \label{eq1}
\begin{split}
\Delta F & = r_a\phi_a^2 + u_a\phi_a^4 + r_c\phi_c^2 + u_c\phi_c^4 + g\phi_a^2\phi_c^2 \\
& + \lambda_a^{xx}\epsilon_{xx}\phi_a^2+\lambda_a^{zz}\epsilon_{zz}\phi_a^2+ \lambda_c^{xx}\epsilon_{xx}\phi_c^2 + \lambda_c^{zz}\epsilon_{zz}\phi_c^2 
\end{split}
\end{equation}
%
where $\phi_a$ and $\phi_c$ describe the amplitude of the two incommensurate CDW states. Here, to account {\it for the fact that the competition between the phases is sufficiently strong to yield a bicritical point, } we consider $g > 0$. The strain-coupling terms have been defined earlier, and we neglect equivalent terms in $\epsilon_{yy}$ for simplicity. In the simplest possible phase diagram as a function of $\epsilon_{xx}, \epsilon_{zz}$ and temperature, the planes defining $T_{CDW}$ for each phase meet along a line of bicritical points.\cite{Pandey2023} The experiment has a fixed relation between $\epsilon_{xx}$ and $\epsilon_{zz}$ defined by the Poisson ratio of Ti (the platform to which the thin sample is bonded) and the differential thermal contraction of Ti and ErTe$_3$, and thus we observe (Figure 4(a)) two lines which meet at a single bicritical point. 
 

At the bicritical point, the coefficients of quadratic terms in the free energy, $r_a$ and $r_c$, have both been tuned to zero. Neglecting for now the strain terms, it is clear that a simple redefinition of the units of either of the CDW order parameters can then lead to an equivalence of the quartic terms, such that exchange of $\phi_a$ and $\phi_c$ leaves the free energy unchanged. In other words, to quartic order in $\phi_a$ and $\phi_c$, and neglecting other terms in the free energy, there is an emergent tetragonal symmetry associated with the critical degrees of freedom (and therefore with all thermodynamic and transport properties that are sensitive to the critical fluctuations). 

However, within this mean field picture, this is clearly only an approximate symmetry. Considering the strain terms, if there were strict tetragonal symmetry (such that $\lambda_c^{xx}=\lambda_a^{zz}=-\lambda_a^{xx}$ etc) then rotation of the x and z axes would leave the free energy invariant. Allowing inequivalence of these terms, however, or including gradient terms or other higher-order terms, all lead to an inequivalence of the $a$ and $c$ crystal directions. Thus, the mean field perspective is one in which it is natural to anticipate an emergent \emph{approximate} higher symmetry close to the bicritical point – the fluctuating order renders the orthorhombic material effectively almost tetragonal, at least with respect to quantities related to the critical fluctuations. To the level of precision of the present experiments, this mean field perspective is sufficient to account for the divergence of the elastoresistivity coefficient $\eta$ upon approach to the putative bicritical point. Remarkably, a thorough treatment of the fluctuations, which we develop in a separate article, reveals that under certain circumstances, this emergent symmetry is in fact asymptotically exact upon approach to criticality \cite{Pandey2023}, though as mentioned above, it is unlikely that experiments with the present control of strain homogeneity could distinguish the associated proposed scaling behavior from the mean field expectations. 

\section{Concluding remarks}

In closing, we comment on some wider implications of our observations. First, having established that modest strains can rotate the direction of the CDW in ErTe$_3$, we note that domain formation is to be anticipated whenever a crystal of RTe$_3$ is held under some conditions that fix the lateral dimensions, at least as long as the spontaneous strain that develops at T$_{CDW}$ is greater than the strain that is caused by fixing the sample dimensions. In particular, one should anticipate domain formation whenever thin samples are bonded to a platform. Since this is a standard geometry for many experiments, consideration of domain formation should be included in interpretation of a variety of experimental results. Indeed, earlier STM results showing domain formation on the surface of TbTe$_3$ crystals might be related to this phenomenology. \cite{Fu2016}

Second, while the symmetry element that renders ErTe$_3$ orthorhombic for all values of externally induced strains is a glide plane, other structural motifs, for example the 1d chains in YBa$_2$Cu$_3$O$_{7-\delta}$, can play a similar role. The present observations provide proof of principle that similar phenomenology of emergent higher-symmetries could, at least in principle, be found in such systems. 

Finally, we note that these observations fall within a wider context of materials for which non-symmetry-enforced degeneracies appear to play a crucial role in establishing emergent properties. Other examples include the possible multicomponent superconducting state in UTe$_2$ (another fundamentally orthorhombic material for which degeneracy of two singlet representations has been proposed \cite{Hayes2021}); possible multicomponent superconducting states in Fe-based superconductors (arising from  ‘accidental’ degeneracy of singlet s and d-wave states in these tetragonal materials \cite{Hirschfeld2011}\cite{Grienko2020}); and proposals for other exotic non-symmetry-enforced multicomponent superconducting states in Sr$_2$RuO$_4$. \cite{Kivelson2020}\cite{Clepkens2021}\cite{Romer2021} 

\section{Acknowledgements}

Work at Stanford (AGS, AP, IRF, SAK, AK, MB) was supported by the Department of Energy, Office of Basic Energy Sciences, under contract DE-AC02-76SF00515. AGS was supported in part by an NSF GRFP and DOE SCGSR award. MB was supported in part by a GLAM Postdoctoral Fellowship. J.J.S. was supported by the National Science Foundation MPS-Ascend Postdoctoral Research Fellowship, Award No. 2138167. The work performed at the Advanced Photon Source was supported by the U.S. Department of Energy (DOE), Office of Science, and Office of Basic Energy Sciences under Contract No. DE-AC02-06CH11357.

\section{Appendix A: Measurement of Strains in the Experiment}

\begin{figure}
\includegraphics[scale=1]{appendix_a.jpg}
\caption{\label{figa} Displacement measured from strain gauge on platform as a function of displacement measured from capacitive sensor shown in blue. At low strains, this function is linear (dashed black line). The strain gauge reading starts to deviate substantially from the capacitive reading at ~7um.} 
\end{figure}

In this experimental set-up, a uniaxial stress is applied to the sample platform producing strains $\Delta L_x/L_x$, $\Delta L_y/L_y$, and $\Delta L_z/L_z$ in the platform. $L_x$ is the effective or active length of the platform which, for a similar shape 'bowtie' has been determined by simulation to be 3.47mm.\cite{Park2020} The capacitive sensor in the strain cell allows for measurement of the displacement between the cell's jaws. For low strains we approximate that this displacement is entirely transmitted to the active (narrow) region of the 'bowtie' platform, giving a measurement of $\Delta L_x$. (In practice though, this transmission is likely between 80-90$\%$.) Once the titanium  plastically deforms, however, (at approximately 0.2$\%$), the transmission of the cell's displacement to the platform falls off significantly, such that the capacitive reading is no longer an accurate determination of $\Delta L_x$. Instead for large displacements, $\Delta L_x$ can be determined by mounting a strain gauge directly onto the Ti platform and a corrective function can be developed to convert the capacitance reading to the actual strain.\ref{figa} This correction is the primary origin of the error bars shown in Fig. 4(a), for example.

While we report our data in terms of measured platform strain, this is a reasonable approximation for the macroscopic strain experienced by the ErTe$3$ sample. Simulations of samples studied reveal that close to 90$\%$ of the platform strain $\Delta L_x/L_x$ is transmitted to the sample. This transmission can also be verified by direct measurement of the sample lattice parameters with XRD. The width of our samples was chosen such that the transverse strain $\Delta L_z/L_z$ is set by the Poisson ratio of the platform rather than the sample. The sample will experience additional strains however due to differential thermal contraction between it and the platform. Since RTe$_3$ contracts 5 times faster than titanium, the sample will experience additional tensile strains as it cools.

The sample's macroscopic strain is differentiated from its microscopic strain in the CDW state, as discussed in the main text. Above the CDW state, $\epsilon_{xx}$ and $\Delta L_x/L_x$ are equivalent. This changes with the presence of domains. While the dimensions of the sample change when the platform is strained, within a single domain $\epsilon_{xx}$ is zero and will not become nonzero until the sample is strained into one of the monodomain states. This can be seen clearly in the Fig. 2b where there is no change in the lattice parameter in the mixed domain phase. 

\section{Appendix B: Determination of T$_{CDW}$}

\begin{figure}
\includegraphics[scale=0.5]{tcdw_plot_paper.png}
\caption{\label{figb} Temperature derivative of $\rho_a +\rho_c$ the longitudinal resistivity, used to determine T$_{CDW}$} 
\end{figure}

As discussed in the text, T$_{CDW}$ is determined by taking a temperature derivative of the longitudinal resistivity and then locating the peak in that derivative as demonstrated in Fig. \ref{figb} . The width of this peak establishes the error. Data sets are offset for clarity. Additionally only a representative subset of the extracted data shown in Figure 4(a) is illustrated here.

\section{Appendix C: Elastoresistivity as a Function of Strain}

\begin{figure*}
  \includegraphics[width=\textwidth,height=6cm]{supplemental_fig2.jpg}
  \caption{\label{figc} Measured values of $\eta$ at fixed values of relative temperature for different definitions of relative temperature as defined in the inset of each panel. Each data set is subtracted by its minimum value to facilitate comparison between temperatures.} 
\end{figure*}

The data presented in Figure 4(c) is influenced by the definition of the relative temperature. In Fig. \ref{figc}, the elastoresistivity data are replotted using (1) a smooth function to model $T_{CDW}$ as a function of $(\Delta L_x/L_x)$, as also done in Figure 4(c), (this is the most accurate method, appropriate even far from the bicritical point); (2) using two straight lines to model $T_{CDW}$ as a function of $(\Delta L_x/L_x)$ (this method is less accurate over the entire strain range, but presumably is still a good description close to the bicritical point); and (3) fixing $T_{CDW}$ as a constant (this is clearly unphysical, but is included to demonstrate how robust the result is). Comparing Figures 7(a) and 7(b), it can be seen that although the exact shape of the peak in the elastoresistance changes  depending on the values of $T_{CDW}$ as a function of $(\Delta L_x/L_x)$, nevertheless the peak feature is robust. If $T_{CDW}$ is assumed to be constant (Figure 7(c)), the peak feature is less pronounced, because the elastoresistivity response grows rapidly at high strains due to the strong enhancement of $T_{CDW}$ which is here neglected. Nevertheless, a peak in the elastoresistivity is still observed close to $T_{CDW}$ even in this case.


%%% OLD DISCUSSION TEXT

%Anticipating then a coupling between charge density wave order and uniaxial strain we now admit terms in our free energy expansion, $\Delta F$, of the form, $\lambda_c^{xx}\epsilon_{xx}\phi_c^2$ where $\phi_c$ is the CDW order parameter which can only enter the free energy expansion in even powers because of the symmetry properties of an incommensurate CDW. Similar terms would be written for $\epsilon_{yy}$ and $\epsilon_{zz}$. There are two immediate consequences of this coupling. First, we expect a spontaneous strain or distortion of the sample as we cool through the CDW transition which indeed has been previously observed.\cite{Ru2008} $Second, T$_{CDW}$ will be tuned by strain.  To complete this free energy expansion though, we must also consider the competing order parameter, CDW order along the “a” direction, $\phi_a$.

%The observed behavior can be understood from a simple mean field picture. Considering both competing CDW order parameters and their coupling the strain the free energy expansion up to quartic order is: $\Delta F = a\phi_c^2 + b\phi_a^2 + d\phi_c^4 + e\phi_a^4 + g\phi_c^2\phi_a^2 +\lambda_c\epsilon_{xx}\phi_c^2 \lambda_a\epsilon_{xx}\phi_a^2$. A phase diagram that can be constructed from this free energy model and is consistent with our experimental findings is shown in Figure 1c and predicts a switching between a and c-axis CDW order with strain.  As discussed, this phase diagram predicts a bicritical point where the phase transitions of both the c-axis and a-axis CDW states coincide. Through this point there will be an identical vanishing of the coefficients of the quadratic order parameter terms, $a\phi_c^2$ and $b\phi_a^2$ in the free energy expansion where $a=a_0(T-T_{C_a}^0)$ and $b=b_0(T-T_{C_c}^0)$. The remaining terms in the free energy expansion can be renormalized such that they are invariant under switching of the a and c axes indicating an emergent tetragonality at the bicritical point. From this mean field perspective, this emergent tetragonality is only approximate since higher order terms have been neglected. However, sufficiently close to T$_{CDW}$, this is a valid approximation. Moreover a full description of the system including fluctuations reveals an exact emergent tetragonality in the critical degrees of freedom.

%The emergence of a higher symmetry via uniaxial strain tuning is non-trivial since the glide plane in the material never is removed and hence the underlying crystal structure of the material remains fundamentally orthorhombic. This framework for strain tuning emergent symmetries could have wider applications in other material systems. In Re$_3$, the restoration of rotational symmetry allows the study of nematic phenomenology in a fundamentally orthorhombic material. In particular, it is possible the material could harbor a vestigial nematic phase, a nematic phase arising from CDW or spin density wave (SDW) fluctuations. The vestigial nematic phase has been proposed as an integral part of both the Fe- based and cuprate superconductors and could play a role in high-T$_c$ physics. The ability to study this phase in the simplified Re$_3$ system would be a valuable tool.

% for the sample file to show the different styles of references, but authors
% most likely will not want to use it.
\nocite{*}

\providecommand{\noopsort}[1]{}\providecommand{\singleletter}[1]{#1}%
\begin{thebibliography}{10}

\bibitem{Monceau2014}
P. Monceau, A.Sinchenko, P.Grigoriev, and P.Lejay.
\newblock Spontaneous breaking of isotropy observed in the electronic transport
  of rare-earth tritellurides.
\newblock {\em Phys. Rev. Lett.}, 112:36601, 2014.

\bibitem{Chu2012}
Jiun-Haw Chu, Hsueh-Hui Kuo, James~G. Analytis, and Ian~R. Fisher.
\newblock Divergent nematic susceptibility in an iron arsenide superconductor.
\newblock {\em Science}, 337, 2012.

\bibitem{Clepkens2021}
Johnathan Clepkens, Austin~W. Lindquist, Xiaoyu Liu, and Hae-Young Kee.
\newblock Higher angular momentum pairings in interorbital shadowed-triplet
  superconductors: Application to sr2ruo4.
\newblock {\em Phys. Rev. B}, 104, 2021.

\bibitem{Eckberg202}
C.~et~al. Eckberg.
\newblock Sixfold enhancement of superconductivity in a tunable electronic
  nematic system.
\newblock {\em Nat. Phys.}, 16, 2020.

\bibitem{Pandey2023}
A.~Pandey et~al.
\newblock Tbd.
\newblock {\em TBD}, 2023.

\bibitem{Fernandes2016}
R~M Fernandes, S~A Kivelson, and E~Berg.
\newblock {Vestigial chiral and charge orders from bidirectional spin-density
  waves: Application to the iron-based superconductors}.
\newblock {\em PHYSICAL REVIEW B}, 93:14511, 2016.

\bibitem{Fu2016}
Ling Fu, Aaron~M. Kraft, Bishnu Sharma, Manoj Singh, Philip Walmsley, Ian~R.
  Fisher, and Michael~C. Boyer.
\newblock Multiple charge density wave states at the surface of tbt e3.
\newblock {\em Physical Review B}, 94, 2016.

\bibitem{Grienko2020}
V.~Grinenko, R.~Sarkar, K.~Kihou, C.~Lee, I.~Morozov, S.~Aswartham, B.~Buchner,
  P.~Chekhonin, W.~Skrotzki, K.~Nenkov, R.~Hu¨hne, K.~Nielsch, S.~L.
  Drechsler, V.~L. Vadimov, M.~A. Silaev, P.~A. Volkov, I.~Eremin, H.~Luetkens,
  and H.~H. Klauss.
\newblock Superconductivity with broken time-reversal symmetry inside a
  superconducting s-wave state.
\newblock {\em Nature Physics}, 16, 2020.

\bibitem{Hayes2021}
I.~M. Hayes, D.~S. Wei, T.~Metz, J.~Zhang, Y.~S. Eo, S.~Ran, S.~R. Saha,
  J.~Collini, N.~P. Butch, D.~F. Agterberg, A.~Kapitulnik, and J.~Paglione.
\newblock Multicomponent superconducting order parameter in ute2.
\newblock {\em Science}, pages 797--801, 2021.

\bibitem{Hicks2014}
Clifford~W Hicks, Mark~E Barber, and Stephen~D Edkins.
\newblock {Piezoelectric-based apparatus for strain tuning}.
\newblock {\em Rev. Sci. Instrum}, 85:65003, 2014.

\bibitem{Hirschfeld2011}
P.~J. Hirschfeld, M.~M. Korshunov, and I.~I. Mazin.
\newblock Gap symmetry and structure of fe-based superconductors.
\newblock {\em Rep. Prog. Phys.}, 74, 2011.

\bibitem{Hristov2019}
Alexander~T. Hristov, Matthias~S. Ikeda, Johanna~C. Palmstrom, Philip Walmsley,
  and Ian~R. Fisher.
\newblock Elastoresistive and elastocaloric anomalies at magnetic and
  electronic-nematic critical points.
\newblock {\em Phys. Rev. B}, 99, 2019.

\bibitem{Hristov2018}
Alexander~T Hristov, Johanna~C Palmstrom, Joshua A~W Straquadine, Tyler~A Merz,
  Harold~Y Hwang, and Ian~R Fisher.
\newblock Measurement of elastoresistivity at finite frequency by amplitude
  demodulation.
\newblock {\em Rev. Sci. Instrum.}, 89:103901, 2018.

\bibitem{Kivelson2020}
Steven~Allan Kivelson, Andrew~Chang Yuan, Brad Ramshaw and Ronny Thomale.
\newblock A proposal for reconciling diverse experiments on the superconducting
  state in sr$_2$ruo$_4$.
\newblock {\em npj Quantum Materials}, 5, 2020.

\bibitem{Kuo2016}
H.~H. Kuo, J.~H. Chu, J.~C. Palmstrom, S.~A. Kivelson, and I.~R. Fisher.
\newblock Ubiquitous signatures of nematic quantum criticality in optimally
  doped fe-based superconductors.
\newblock {\em Science}, 352, 2016.

\bibitem{Kuo2013}
Hsueh-Hui Kuo, Maxwell~C. Shapiro, Scott~C. Riggs, and Ian~R. Fisher.
\newblock Measurement of the elastoresistivity coefficients of the underdoped
  iron arsenide ba(fe0.975co0.025)2as2.
\newblock {\em Phys. Rev. B}, 88, 2013.

\bibitem{Sanchez2019}
Paul Malinowski, Qianni Jiang, Joshua Sanchez, Zhaoyu Liu, Joshua Mutch,
  Preston Went, Jian Liu, Philip Ryan, Jong-Woo Kim, and Jiun-Haw Chu.
\newblock {Drastic suppression of superconducting {\$}T{\_}{\{}c{\}}{\$} by
  anisotropic strain near a nematic quantum critical point}.
\newblock {\em Nature Physics}, pages 1--5, nov 2019.

\bibitem{Maschek2015}
M~Maschek, S.~Rosenkranz, R.~Heid, A.~H. Said, P.~Giraldo-Gallo, I.~R. Fisher,
  and F.~Weber.
\newblock {Wave-Vector-Dependent Electron-Phonon Coupling and the
  Charge-Density-Wave Transition in TbTe3}.
\newblock {\em Phys. Rev. B}, 91, 2015.

\bibitem{Maschek2018}
M.~Maschek, D.~A. Zocco, S.~Rosenkranz, R.~Heid, A.~H. Said, A.~Alatas,
  P.~Walmsley, I.~R. Fisher, and F.~Weber.
\newblock {Competing soft phonon modes at the charge-density-wave transitions
  in DyT e3}.
\newblock {\em Physical Review B}, 98(9), 2018.

\bibitem{Nie2017}
Laimei Nie, Akash~V Maharaj, Eduardo Fradkin, and Steven~A Kivelson.
\newblock {Vestigial nematicity from spin and/or charge order in the cuprates}.
\newblock {\em Phys. Rev. B}, 96(085142), 2017.

\bibitem{Nie2014}
Laimei Nie, Gilles Tarjus, and Steven~Allan Kivelson.
\newblock {Quenched disorder and vestigial nematicity in the pseudogap regime
  of the cuprates}.
\newblock {\em Proceedings of the National Academy of Sciences of the United
  States of America}, 111(22):7980--7985, jun 2014.

\bibitem{Park2020}
Joonbum Park, Jack~M Bartlett, Hilary M~L Noad, Alexander~L Stern, Mark~E
  Barber, Markus K{\"{o}}nig, Suguru Hosoi, Takasada Shibauchi, Andrew~P
  Mackenzie, Alexander Steppke, and Clifford~W Hicks.
\newblock {Rigid platform for applying large tunable strains to mechanically
  delicate samples}.
\newblock {\em Review of Scientific Instruments}, 91, 2020.

\bibitem{Romer2021}
Astrid Romer, PJ~Hirschfeld, and Brian~M. Andersen.
\newblock Superconducting state of sr2ruo4 in the presence of longer-range
  coulomb interactions.
\newblock {\em Phys. Rev. B}, 104, 2021.

\bibitem{Ru2008}
N~Ru, C~L Condron, G~Y Margulis, K~Y Shin, J~Laverock, S~B Dugdale, M~F Toney,
  and I~R Fisher.
\newblock {Effect of chemical pressure on the charge density wave transition in
  rare-earth tritellurides RTe 3}.
\newblock {\em Phys. Rev. B}, 77:035114, 2008.

\bibitem{Ru2006}
N.~Ru and I.~R. Fisher.
\newblock { Thermodynamic and transport properties of YTe3 , LaTe3 , and CeTe3
  }.
\newblock {\em Physical Review B}, 73(3):3--6, 2006.

\bibitem{Ru_th}
Nancy Ru.
\newblock {\em {Charge Denisty Wave Formation in Rare-Earth Tritellurides}}.
\newblock PhD thesis, Leland Stanford Jr. University, 2008.

\bibitem{Shapiro2015}
M.~C. Shapiro, Patrik Hlobil, A.~T. Hristov, Akash~V. Maharaj, and I.~R.
  Fisher.
\newblock Symmetry constraints on the elastoresistivity tensor.
\newblock {\em Phys. Rev. B}, 92, 2015.

\bibitem{Shapiro2016}
M.~C. Shapiro, A.~T. Hristov, J.~C. Palmstrom, Jiun~Haw Chu, and I.~R. Fisher.
\newblock {Measurement of the B 1 g and B 2 g components of the
  elastoresistivity tensor for tetragonal materials via transverse resistivity
  configurations}.
\newblock {\em Review of Scientific Instruments}, 87(6), 2016.

\bibitem{Straquadine2020}
J~A~W Straquadine, M~S Ikeda, and I~R Fisher.
\newblock {Evidence for realignment of the charge density wave state in ErTe 3
  and TmTe 3 under uniaxial stress via elastocaloric and elastoresistivity
  measurements}.
\newblock {\em arXiv}, 2005.10461, 2020.

\bibitem{Walmsley2017}
P.~Walmsley and I.~R. Fisher.
\newblock {Determination of the resistivity anisotropy of orthorhombic
  materials via transverse resistivity measurements}.
\newblock {\em Review of Scientific Instruments}, 88(4):1--10, 2017.

\bibitem{Yao2006}
Hong Yao, John~A. Robertson, Eun~Ah Kim, and Steven~A. Kivelson.
\newblock Theory of stripes in quasi-two-dimensional rare-earth tellurides.
\newblock {\em Physical Review B - Condensed Matter and Materials Physics}, 74,
  2006.

\end{thebibliography}

\end{document}
Compile
\end{document}
%
% ****** End of file apssamp.tex ******
